\chapter{Server Actions, Email \& Notifications}
\label{ch:server-actions}

\section{What a Server Action is, and how OLK uses them}

A file starting with the string \code{'use server'} exports functions that Next.js turns into callable, server-only RPCs from Client Components — the function body only ever executes on Vercel's servers, but a Client Component can \code{await} it as if it were a local async function, with no hand-written \code{fetch} or API route needed. \pathname{src/lib/admin-actions.ts} and \pathname{src/lib/notifications.ts} are the only two \code{'use server'} files in this codebase; every admin page in Chapter~\ref{ch:admin} calls into the former directly.

\section{The admin action pattern: guard, act, log, revalidate}

Every one of the 24 functions in \pathname{admin-actions.ts} follows the identical four-step shape — but as of a code-quality pass across the whole codebase, that shape now lives in exactly one place, a higher-order function named \code{withAdminAction}, instead of being copy-pasted into every function body:

\begin{lstlisting}[caption={src/lib/admin-actions.ts — the shared wrapper, illustrated with banUser}]
function withAdminAction<Args extends unknown[]>(
  minRole: AdminRole,
  action: string,
  handler: (ctx: AdminActionCtx, ...args: Args) => Promise<AuditMeta>
): (...args: Args) => Promise<ActionResult> {
  return async (...args: Args) => {
    try {
      const admin = await guardRole(minRole)                    // 1. GUARD
      const supabase = await createClient()
      const { targetType, targetId, details } =
        await handler({ admin, supabase }, ...args)              // 2. ACT
      await auditLog(admin.id, action, targetType, targetId, details)  // 3. LOG
      revalidatePath('/admin')                                   // 4. REVALIDATE
      return { success: true }
    } catch (e) {
      return { success: false, error: (e as Error).message }
    }
  }
}

export const banUser = withAdminAction(
  'moderator', 'ban_user',
  async ({ admin, supabase }, userId: string, reason: string, isPermanent: boolean, durationDays?: number) => {
    const expiresAt = !isPermanent && durationDays
      ? new Date(Date.now() + durationDays * 86400000).toISOString() : null

    await supabase.from('user_bans').insert({ user_id: userId, admin_id: admin.id, reason, is_permanent: isPermanent, expires_at: expiresAt })
    await supabase.from('profiles').update({ is_banned: true, ban_reason: reason, ban_expires_at: expiresAt }).eq('id', userId)
    await supabase.from('notifications').insert({ user_id: userId, type: 'admin', title: /* ... */ '', link: '/profile' })

    return { targetType: 'user', targetId: userId, details: { reason, isPermanent, durationDays } }
  }
)
\end{lstlisting}

\begin{enumerate}
  \item \textbf{Guard} — \code{guardRole(minRole)} calls \code{checkAdminAPI} (Chapter~\ref{ch:rls}) and throws if the caller doesn't meet the role floor. This is a check made on top of RLS, not instead of it — the subsequent Supabase calls in the handler still go through RLS as whatever role the Server Action's own Supabase client authenticates as (the calling user's session, via \pathname{src/lib/supabase/server.ts}), so a bug in \code{guardRole} would still have to also defeat RLS to cause real damage.
  \item \textbf{Act} — the handler passed to \code{withAdminAction} does the actual mutation(s), near-always inserting a \code{notifications} row telling the affected user what just happened to them, then returns \emph{what to audit-log} rather than performing the log itself.
  \item \textbf{Log} — the wrapper writes exactly one row to \code{admin\_audit\_log} (Chapter~\ref{ch:schema}) per call, built from the handler's return value plus the wrapper's own \code{action} string. This is not optional or conditional anywhere — a handler cannot forget it, because it isn't the handler's job to call it.
  \item \textbf{Revalidate} — \code{revalidatePath('/admin')} tells Next.js's cache that server-rendered data under \pathname{/admin} may be stale and should be refetched on next navigation.
\end{enumerate}

\begin{olktip}
Every action still resolves to \code{\{ success: boolean, error?: string \}} rather than throwing across the server/client boundary — that contract is unchanged, only where the \code{try/catch} lives has moved. Calling code throughout the admin pages (Chapter~\ref{ch:pages}) still uniformly checks \code{result.success} rather than using \code{try/catch} at the call site.
\end{olktip}

\begin{olkhowto}
Adding action \#25 to \pathname{admin-actions.ts}? Wrap it with \code{withAdminAction} rather than hand-writing the guard/log/revalidate scaffolding:
\begin{enumerate}
  \item \code{export const yourAction = withAdminAction(minRole, 'your\_action\_name', async (\{ admin, supabase \}, ...yourArgs) => \{ ... \})} — \code{minRole} is \code{'moderator'} for anything reversible, \code{'super\_admin'} for anything destructive or platform-settings-adjacent (Chapter~\ref{ch:admin}).
  \item Inside the handler: mutate, then insert a \code{notifications} row for whichever user should be told what just happened.
  \item End the handler with \code{return \{ targetType: 'your\_target\_type', targetId, details: \{ ...whatever's worth auditing \} \}} — this is what becomes the \code{admin\_audit\_log} row; there is no way to skip it, since it's also the handler's only return value.
\end{enumerate}
The old failure mode — an action that works but is invisible to \pathname{admin/settings/page.tsx}'s audit log because a developer forgot the \code{auditLog(...)} call — is now structurally harder to hit, since logging happens in the wrapper rather than being one more line a handler author has to remember.
\end{olkhowto}

\section{Notifications: in-app row first, email best-effort second}

\begin{lstlisting}[caption={src/lib/notifications.ts — in full}]
type NotificationContext =
  | { kind: 'request'; id: string }
  | { kind: 'club_join'; id: string }
  | { kind: 'club_announcement'; id: string }
  | { kind: 'event_created'; id: string }
  | { kind: 'wishlist_match'; id: string }

export async function createNotification({ userId, type, title, link, context }: NotificationPayload) {
  const supabase = await createServerClient()          // the calling user's session
  const { data: { user } } = await supabase.auth.getUser()
  if (!user) throw new Error('Unauthorized')

  const { data: allowed, error: guardError } = await supabase.rpc('can_notify', {
    p_target_user: userId,
    p_context_type: context.kind,
    p_context_id: context.id,
  })
  if (guardError || !allowed) {
    throw new Error('Forbidden: no valid relationship to notify this user')
  }

  const supabaseAdmin = createAdminClient(              // service-role client
    process.env.NEXT_PUBLIC_SUPABASE_URL!,
    process.env.SUPABASE_SERVICE_ROLE_KEY!,
  )
  const { error } = await supabaseAdmin.from('notifications').insert({
    user_id: userId, type, title, link: link || '/notifications',
  })
  if (error) throw error

  try {
    await sendNotificationEmail({ userId, type, title })
  } catch (err) {
    console.error('Email notification failed:', err)
  }
}
\end{lstlisting}

The in-app notification row is written first and its failure is a hard error (thrown, propagated to the caller); the email send is wrapped in its own \code{try/catch} and its failure is only logged. This ordering is deliberate: a user should always see the notification inside the app even if their email is undeliverable or Resend is down — email is a convenience layer on top of a system that must work without it.

\begin{olknote}
This function used to insert with the calling user's own session client, which worked only because the \code{notifications} INSERT policy was, at the time, wide open to any authenticated user (Chapter~\ref{ch:rls}). Now that policy only allows \code{auth.uid() = user\_id}, and most calls here are inherently cross-user (the \emph{other} party in a request/message/club gets notified, not the caller) — so this action switched to the service-role client, after first confirming the caller has a valid session. It is now the trusted server-side gate for every cross-user notification the app creates outside the admin actions in Chapter~\ref{ch:admin} (which have their own, separate "admins/mods can insert for anyone" policy).
\end{olknote}

\begin{olknote}
\textbf{That gate used to check only "is the caller logged in," not "does the caller actually have a reason to notify this specific person."} Since every call here goes through the service-role client, an authenticated-only check meant any user could call this Server Action directly (it's a plain exported async function, reachable from any client) with an arbitrary \code{userId}, \code{title}, and \code{link} — spamming or phishing an unrelated user with a fake notification. The fix is the \code{context} parameter and the \code{can\_notify(p\_target\_user, p\_context\_type, p\_context\_id)} RPC shown above: a \code{SECURITY DEFINER} SQL function (Chapter~\ref{ch:functions}) that verifies the caller and \code{userId} are the two real parties on the referenced \code{book\_requests}/\code{clubs}/\code{club\_members}/\code{wishlists} row, before the insert is allowed to proceed. It needs to be \code{SECURITY DEFINER} for at least one case that RLS alone can't express: a book owner adding a new listing that matches someone else's wishlist has no RLS-visible relationship to that wishlist row (\code{"Users can view own wishlists"} would otherwise hide it from them entirely), so the verification has to run with elevated privilege and check the real underlying fact — that the wishlist's \code{matched\_book\_id} now points to a book the caller owns — rather than relying on what the caller's own session is allowed to \code{SELECT}. All six call sites (\pathname{browse}, \pathname{my-books}, \pathname{requests}, \pathname{messages/[id]}, \pathname{clubs/[id]} $\times$2) now pass a \code{context} identifying which real-world relationship justifies the notification.
\end{olknote}

\section{Sending the actual email: a service-role client, on purpose}

\begin{lstlisting}[caption={src/lib/send-notification-email.ts — the privileged lookup}]
export async function sendNotificationEmail({ userId, type, title }: {...}) {
  if (!resend) return   // silently no-op if RESEND_API_KEY is unset

  const supabaseAdmin = createClient(
    process.env.NEXT_PUBLIC_SUPABASE_URL!,
    process.env.SUPABASE_SERVICE_ROLE_KEY!,   // NOT the anon key
  )
  const { data: authUser } = await supabaseAdmin.auth.admin.getUserById(userId)
  const email = authUser?.user?.email
  if (!email) return
  // ... resend.emails.send({ ..., html: `... ${escapeHtml(title)} ...` })
}
\end{lstlisting}

Resolving \emph{any} user's email address by id is an operation no ordinary RLS policy would ever grant a client — email addresses live in \code{auth.users}, not \code{public.profiles}, and are never exposed through the query builder. This file constructs a second, separate Supabase client using the \textbf{service role key} instead of the anon key specifically to call the privileged \code{auth.admin.getUserById} API, and the file begins with \code{import 'server-only'} (Chapter~\ref{ch:techstack}) so that a build error occurs if this file is ever accidentally reachable from client-side code — the service role key must never reach a browser bundle.

\begin{olkwarn}
Notice \code{escapeHtml(title)} around the one piece of user-influenced content interpolated into the email's raw HTML string. A notification's \code{title} field can, depending on the notification type, contain text a regular user chose (e.g. a report reason echoed back, or in principle a maliciously crafted book title flowing through some future notification type). Escaping it before string-interpolating into HTML is the only thing standing between that and an HTML/script injection inside an email client. If you add a new notification type, keep any user-influenced text passed through \code{escapeHtml} before it reaches this file.
\end{olkwarn}

\begin{olkhowto}
A brand-new notification \code{type} (say \code{'club\_invite'}) touches more than the string you pass to \code{createNotification}:
\begin{enumerate}
  \item There is no enum or TypeScript union constraining \code{type} — it's a plain string column (Chapter~\ref{ch:schema}). A typo here doesn't fail at write time; it fails silently at read time, as step 2 below.
  \item \pathname{src/components/notification-bell.tsx} and \pathname{src/app/(dashboard)/notifications/page.tsx} both index into an \code{ICONS} lookup object keyed by \code{n.type}, falling back to a generic bell emoji when the key is missing — add your new type as a key in \emph{each} file's \code{ICONS} object, or it silently falls back to that generic icon rather than erroring.
  \item If \code{title} ever includes user-typed text for this new type, route it through \code{escapeHtml} before it can reach \pathname{send-notification-email.ts} (the warning above) — easy to forget for a brand-new call site since nothing enforces it.
  \item Decide which existing \code{context.kind} the new call site's relationship maps to (\code{'request'}, \code{'club\_join'}, \code{'club\_announcement'}, \code{'event\_created'}, or \code{'wishlist\_match'}) — or, if it's a genuinely new kind of relationship, add a new branch to \code{can\_notify()} (Chapter~\ref{ch:rls}) \emph{before} wiring up the call site. \code{'event\_created'} itself is the \code{'club\_announcement'} shape one table over: the caller must be the event's own \code{creator\_id}, and the target must be a real member of that event's club. \code{createNotification} throws \code{"Forbidden"} if \code{can\_notify} doesn't recognise the \code{context.kind} or can't verify the relationship, so a new notification type with no matching branch fails loudly for every caller rather than silently succeeding.
\end{enumerate}
\end{olkhowto}

\section{The weekly digest: a genuine Route Handler}

\pathname{src/app/api/digest/route.ts} is the one place in this codebase that is a real HTTP endpoint rather than a Server Action, because it needs to be invoked by something outside a page render — an external cron scheduler (Chapter~\ref{ch:vercel}).

\begin{lstlisting}[caption={app/api/digest/route.ts — auth gate, batched email lookup, batched send}]
export const maxDuration = 60           // Vercel function timeout budget
const RESEND_BATCH_SIZE = 100           // Resend's batch-send API limit

async function handleDigest(request: Request) {
  const authHeader = request.headers.get('authorization')
  if (!isAuthorized(authHeader)) {      // timing-safe comparison against CRON_SECRET
    return Response.json({ error: 'Unauthorized' }, { status: 401 })
  }
  if (!resend) return Response.json({ error: 'Email service not configured' }, { status: 503 })

  const supabaseAdmin = createClient(process.env.NEXT_PUBLIC_SUPABASE_URL!, process.env.SUPABASE_SERVICE_ROLE_KEY!)
  const subscribers = /* profiles where email_digest = true */
  const newBooks = /* books from the last 7 days, status = available, limit 10 */

  // One id -> email map via paginated listUsers(), not one getUserById() call
  // per subscriber -- O(subscribers / 1000) Auth API calls instead of O(subscribers).
  const emailById = new Map<string, string>()
  for (let page = 1; ; page++) {
    const { data, error } = await supabaseAdmin.auth.admin.listUsers({ page, perPage: 1000 })
    if (error || !data?.users?.length) break
    for (const u of data.users) if (u.email) emailById.set(u.id, u.email)
    if (data.users.length < 1000) break
  }

  const recipients = subscribers
    .map(sub => ({ sub, email: emailById.get(sub.id) }))
    .filter((r): r is { sub: typeof subscribers[number]; email: string } => !!r.email)

  let sent = 0
  const failures: { userIds: string[]; error: string }[] = []
  for (const batch of chunk(recipients, RESEND_BATCH_SIZE)) {
    try {
      await resend.batch.send(batch.map(({ sub, email }) => ({ from: fromAddress, to: email, subject, html: renderEmail(sub) })))
      sent += batch.length
    } catch (e) {
      failures.push({ userIds: batch.map(r => r.sub.id), error: (e as Error).message })
    }
  }
  return Response.json({ sent, books: newBooks.length, failures: failures.length, failureDetails: failures })
}

export async function GET(request: Request) { return handleDigest(request) }
export async function POST(request: Request) { return handleDigest(request) }
\end{lstlisting}

Authorization here is a timing-safe shared-secret Bearer token comparison against \code{CRON\_SECRET} (Chapter~\ref{ch:vercel} covers where that value comes from) — appropriate for a server-to-server cron trigger with no user session involved.

\begin{olknote}
The per-subscriber loop used to have no \code{try/catch} around \code{resend.emails.send(...)} — a malformed address or a transient Resend error would abort the entire run for every subscriber after that point in iteration order. It also used to call \code{auth.admin.getUserById(sub.id)} and \code{resend.emails.send(...)} once each, per subscriber, sequentially — a real N+1 that would exceed Vercel's function timeout well before this project's 30,000--50,000-user target (Chapter~\ref{ch:welcome}). The 23-07-2026 scaling pass replaced both: email lookups now go through paginated \code{listUsers()} building one id$\to$email map, and sends go through \code{resend.batch.send()} in chunks of 100 rather than one \code{resend.emails.send()} call per person. Failures are still collected per-batch rather than aborting the run, and an explicit \code{maxDuration} now states the timeout budget rather than relying on Vercel's default. The route also used to export only \code{POST}, while Vercel's Cron Jobs feature calls with \code{GET} by default (Chapter~\ref{ch:vercel}); it now exports both, delegating to the same \code{handleDigest} logic.
\end{olknote}

\begin{olktask}
With \code{RESEND\_API\_KEY} unset (the default in local development, per Chapter~\ref{ch:local-env}), trigger the digest route locally with \code{curl -X POST localhost:3000/api/digest -H "Authorization: Bearer \$CRON\_SECRET"} and confirm you get the \code{503} "Email service not configured" response. Then set a (fake, non-functional) \code{RESEND\_API\_KEY} in \pathname{.env.local}, restart the dev server, and confirm the response changes to an attempted send that now fails differently — this is the fastest way to see the file's two independent failure modes without needing a real Resend account.
\end{olktask}
